% GENERATED FILE. Edit canonical lessons, then run npm run generate:books.
% canonical-source-hash: fnv1a64:b4a8d56e6da84125
% canonical-lessons: BN-C296-songbadpotro, BN-C296-shironam, BN-C296-protibedon, BN-C296-prokash-kora, BN-C296-sompadok

\chapter{The Morning News}
\label{ch:bn-the-morning-news}
\hlchaptermodality{296}

\begin{chapteropening}
\textbf{By the end of this chapter:} \emph{I can talk about newspapers, headlines and reports.}

Complete the last lesson of chapter 296: I can talk about newspapers, headlines and reports.
\end{chapteropening}

\section[\texorpdfstring{\bn{সংবাদপত্র} (sôngbādpôtrô)}{সংবাদপত্র (sôngbādpôtrô)}]{\bn{সংবাদপত্র} (sôngbādpôtrô) — a newspaper}

\label{lesson:BN-C296-songbadpotro}



\begin{quote}
Before the new one: say the Bengali for to tackle, then the Bengali for to regain.
\end{quote}

\subsection*{You'll want to know: \bn{সংবাদপত্র}}
\textbf{\bn{সংবাদপত্র}} — \emph{sôngbādpôtrô} — \textquotedblleft{}a newspaper\textquotedblright{}.

\begin{quote}
\bn{বাবা} \bn{প্রতিদিন} \bn{সকালে} \bn{সংবাদপত্র} \bn{পড়েন।}
\end{quote}

\begin{quote}
\emph{bābā prôtidin sôkāle sôngbādpôtrô pôṛen.}
\end{quote}

\begin{quote}
Father reads the newspaper every morning.
\end{quote}

\begin{cousinweb}
\textbf{\bn{সংবাদপত্র}} (\emph{sôngbādpôtrô}) joins Sanskrit \emph{saṃvāda}, 'a speaking together, a conversation' (Monier-Williams), to \emph{pattra}, 'a leaf, a written sheet': a sheet of news. Everyday speech says \bn{খবরের} \bn{কাগজ}.
\end{cousinweb}

\begin{grammarlens}[title={What you've built}]
One more word for telling what happened.
\end{grammarlens}

\subsection*{Your turn}
\begin{itemize}
\raggedright
  \item \emph{Say it:} \emph{sôngbādpôtrô}
  \item \emph{Say it:} \emph{sôngbādpôtrô}, once more
  \item \emph{Say it:} \emph{bābā prôtidin sôkāle sôngbādpôtrô pôṛen.} — Father reads the newspaper every morning.
  \item \emph{Recall:} say the Bengali for to tackle, then the Bengali for to regain, then say \emph{sôngbādpôtrô} again
\end{itemize}

\begin{tcolorbox}[breakable,colback=teal!4,colframe=teal!35!black,title={Before you move on}]
What does \bn{সংবাদপত্র} mean? (\textquotedblleft{}A newspaper\textquotedblright{}.) Say it once more.
\end{tcolorbox}
\section[\texorpdfstring{\bn{শিরোনাম} (shironām)}{শিরোনাম (shironām)}]{\bn{শিরোনাম} (shironām) — a headline}

\label{lesson:BN-C296-shironam}



\begin{quote}
Before the new one: say the Bengali for to regain, then the Bengali for a newspaper.
\end{quote}

\subsection*{You'll want to know: \bn{শিরোনাম}}
\textbf{\bn{শিরোনাম}} — \emph{shironām} — \textquotedblleft{}a headline\textquotedblright{}.

\begin{quote}
\bn{আজ} \bn{সংবাদপত্রের} \bn{বড়} \bn{শিরোনাম} \bn{ছিল} \bn{ঘূর্ণিঝড়ের} \bn{খবর।}
\end{quote}

\begin{quote}
\emph{āj sôngbādpôtrer bôṛo shironām chhilô ghūrnijhôṛer khôbôr.}
\end{quote}

\begin{quote}
Today the newspaper's big headline was the news of the cyclone.
\end{quote}

\begin{cousinweb}
\textbf{\bn{শিরোনাম}} (\emph{shironām}) joins Sanskrit \emph{śiras}, 'the head, the top' (Monier-Williams), to \textbf{\bn{নাম}}, a name: the name at the head of a story. It is a modern learned compound.
\end{cousinweb}

\begin{grammarlens}[title={What you've built}]
One more word for telling what happened.
\end{grammarlens}

\subsection*{Your turn}
\begin{itemize}
\raggedright
  \item \emph{Say it:} \emph{shironām}
  \item \emph{Say it:} \emph{shironām}, once more
  \item \emph{Say it:} \emph{āj sôngbādpôtrer bôṛo shironām chhilô ghūrnijhôṛer khôbôr.} — Today the newspaper's big headline was the news of the cyclone.
  \item \emph{Recall:} say the Bengali for to regain, then the Bengali for a newspaper, then say \emph{shironām} again
\end{itemize}

\begin{tcolorbox}[breakable,colback=teal!4,colframe=teal!35!black,title={Before you move on}]
What does \bn{শিরোনাম} mean? (\textquotedblleft{}A headline\textquotedblright{}.) Say it once more.
\end{tcolorbox}
\section[\texorpdfstring{\bn{প্রতিবেদন} (prôtibedôn)}{প্রতিবেদন (prôtibedôn)}]{\bn{প্রতিবেদন} (prôtibedôn) — a report}

\label{lesson:BN-C296-protibedon}



\begin{quote}
Before the new one: say the Bengali for a newspaper, then the Bengali for a headline.
\end{quote}

\subsection*{You'll want to know: \bn{প্রতিবেদন}}
\textbf{\bn{প্রতিবেদন}} — \emph{prôtibedôn} — \textquotedblleft{}a report\textquotedblright{}.

\begin{quote}
\bn{সাংবাদিক} \bn{বন্যার} \bn{ওপর} \bn{একটা} \bn{লম্বা} \bn{প্রতিবেদন} \bn{লিখলেন।}
\end{quote}

\begin{quote}
\emph{sāngbādik bônyār opôr ekṭā lômbā prôtibedôn likhlen.}
\end{quote}

\begin{quote}
The journalist wrote a long report on the flood.
\end{quote}

\begin{cousinweb}
\textbf{\bn{প্রতিবেদন}} (\emph{prôtibedôn}) is a learned word built from Sanskrit \emph{prati-}, 'back, in return', and \emph{vedana}, 'making known' (Monier-Williams): an account carried back to the reader.
\end{cousinweb}

\begin{grammarlens}[title={What you've built}]
One more word for telling what happened.
\end{grammarlens}

\subsection*{Your turn}
\begin{itemize}
\raggedright
  \item \emph{Say it:} \emph{prôtibedôn}
  \item \emph{Say it:} \emph{prôtibedôn}, once more
  \item \emph{Say it:} \emph{sāngbādik bônyār opôr ekṭā lômbā prôtibedôn likhlen.} — The journalist wrote a long report on the flood.
  \item \emph{Recall:} say the Bengali for a newspaper, then the Bengali for a headline, then say \emph{prôtibedôn} again
\end{itemize}

\begin{tcolorbox}[breakable,colback=teal!4,colframe=teal!35!black,title={Before you move on}]
What does \bn{প্রতিবেদন} mean? (\textquotedblleft{}A report\textquotedblright{}.) Say it once more.
\end{tcolorbox}
\section[\texorpdfstring{\bn{প্রকাশ} \bn{করা} (prôkāsh kôrā)}{প্রকাশ করা (prôkāsh kôrā)}]{\bn{প্রকাশ} \bn{করা} (prôkāsh kôrā) — to publish; to make public}

\label{lesson:BN-C296-prokash-kora}



\begin{quote}
Before the new one: say the Bengali for a headline, then the Bengali for a report.
\end{quote}

\subsection*{You'll want to know: \bn{প্রকাশ} \bn{করা}}
\textbf{\bn{প্রকাশ} \bn{করা}} — \emph{prôkāsh kôrā} — \textquotedblleft{}to publish; to make public\textquotedblright{}.

\begin{quote}
\bn{পত্রিকাটা} \bn{ঘটনার} \bn{ছবি} \bn{প্রকাশ} \bn{করল।}
\end{quote}

\begin{quote}
\emph{pôtrikāṭā ghôṭônār chhôbi prôkāsh kôrlô.}
\end{quote}

\begin{quote}
The magazine published pictures of the incident.
\end{quote}

Past (he/she …): \bn{তিনি} \bn{খবরটা} \bn{প্রকাশ} \bn{করলেন}, \emph{tini khôbôrṭā prôkāsh kôrlen}.

\begin{cousinweb}
\textbf{\bn{প্রকাশ} \bn{করা}} (\emph{prôkāsh kôrā}) puts Sanskrit \emph{prakāśa}, 'shining, visible; open, public' (Monier-Williams), before \bn{করা}: to bring something into the light.
\end{cousinweb}

\begin{grammarlens}[title={What you've built}]
One more word for telling what happened.
\end{grammarlens}

\subsection*{Your turn}
\begin{itemize}
\raggedright
  \item \emph{Say it:} \emph{prôkāsh kôrā}
  \item \emph{Say it:} \emph{prôkāsh kôrā}, once more
  \item \emph{Say it:} \emph{pôtrikāṭā ghôṭônār chhôbi prôkāsh kôrlô.} — The magazine published pictures of the incident.
  \item \emph{Recall:} say the Bengali for a headline, then the Bengali for a report, then say \emph{prôkāsh kôrā} again
\end{itemize}

\begin{tcolorbox}[breakable,colback=teal!4,colframe=teal!35!black,title={Before you move on}]
What does \bn{প্রকাশ} \bn{করা} mean? (\textquotedblleft{}To publish; to make public\textquotedblright{}.) Say it once more.
\end{tcolorbox}
\section[\texorpdfstring{\bn{সম্পাদক} (sômpādôk)}{সম্পাদক (sômpādôk)}]{\bn{সম্পাদক} (sômpādôk) — an editor}

\label{lesson:BN-C296-sompadok}



\begin{quote}
Before the new one: say the Bengali for a report, then the Bengali for to publish.
\end{quote}

\subsection*{You'll want to know: \bn{সম্পাদক}}
\textbf{\bn{সম্পাদক}} — \emph{sômpādôk} — \textquotedblleft{}an editor\textquotedblright{}.

\begin{quote}
\bn{সম্পাদক} \bn{প্রতিবেদনটা} \bn{পড়ে} \bn{শিরোনাম} \bn{বদলালেন।}
\end{quote}

\begin{quote}
\emph{sômpādôk prôtibedônṭā pôṛe shironām bôdlālen.}
\end{quote}

\begin{quote}
The editor read the report and changed the headline.
\end{quote}

\begin{cousinweb}
\textbf{\bn{সম্পাদক}} (\emph{sômpādôk}) is Sanskrit \emph{sampādaka}, 'procuring, bringing about' (Monier-Williams); the sense 'an editor', the one who gets a paper made, is the modern one.
\end{cousinweb}

\begin{grammarlens}[title={What you've built}]
One more word for telling what happened.
\end{grammarlens}

\subsection*{Your turn}
\begin{itemize}
\raggedright
  \item \emph{Say it:} \emph{sômpādôk}
  \item \emph{Say it:} \emph{sômpādôk}, once more
  \item \emph{Say it:} \emph{sômpādôk prôtibedônṭā pôṛe shironām bôdlālen.} — The editor read the report and changed the headline.
  \item \emph{Recall:} say the Bengali for a report, then the Bengali for to publish, then say \emph{sômpādôk} again
\end{itemize}

\begin{tcolorbox}[breakable,colback=teal!4,colframe=teal!35!black,title={Before you move on}]
What does \bn{সম্পাদক} mean? (\textquotedblleft{}An editor\textquotedblright{}.) Say it once more.
\end{tcolorbox}
